% Shared preamble for the Korean and English RamFlow papers.
% Each main file defines \ifko (true for Korean) before \input-ing this file.
\usepackage[a4paper,top=24mm,bottom=26mm,left=25mm,right=25mm,headsep=8mm,footskip=11mm]{geometry}
\usepackage{amsmath,amssymb}
\usepackage{mathpazo}
\usepackage[no-math]{fontspec}
\ifko
  % Korean: Noto CJK KR carries both Hangul and Latin glyphs, so one family
  % keeps the colour of the page even. Lines break at inter-word spaces.
  \setmainfont{Noto Serif CJK KR}[Scale=0.94,BoldFont={* SemiBold}]
  \setsansfont{Noto Sans CJK KR}[Scale=0.92,BoldFont={* Bold}]
  \setmonofont{DejaVu Sans Mono}[Scale=0.80]
  \linespread{1.30}
  \XeTeXlinebreaklocale "ko"
  \XeTeXlinebreakskip=0pt plus 0.6pt
  \AtBeginDocument{\renewcommand{\emph}[1]{\textsf{#1}}}
  % inside diagrams and tables, break only at spaces and never letter-space
  \AtBeginDocument{\tikzset{every picture/.append style={execute at begin picture={\XeTeXlinebreaklocale ""\XeTeXlinebreakskip=0pt\relax},execute at end picture={\XeTeXlinebreaklocale "ko"}}}}
  \renewcommand{\figurename}{그림}
  \renewcommand{\tablename}{표}
  \renewcommand{\contentsname}{차례}
  \renewcommand{\appendixname}{부록}
  \renewcommand{\abstractname}{초록}
  \AtBeginDocument{\floatname{algorithm}{알고리즘}}
  \AtBeginDocument{\renewcommand{\figurename}{그림}\renewcommand{\tablename}{표}}
\else
  \setmainfont{TeX Gyre Pagella}
  \setsansfont{Noto Sans CJK KR}[Scale=0.92]
  \setmonofont{DejaVu Sans Mono}[Scale=0.80]
  \linespread{1.12}
\fi
\usepackage{microtype}
\setlength{\emergencystretch}{1.5em}
\usepackage{xcolor}
\definecolor{accent}{HTML}{2A78D6}
\definecolor{offc}{HTML}{EB6834}
\definecolor{ink2}{HTML}{52514E}
\definecolor{rule}{HTML}{D9D8D3}
\definecolor{panel}{HTML}{F4F3F0}
\usepackage{graphicx}
\usepackage{booktabs,tabularx,multirow,array}
\usepackage{caption}
\captionsetup{font=small,labelfont={bf,sf},labelsep=period,skip=5pt}
\usepackage{enumitem}
\setlist{itemsep=1pt,topsep=3pt,leftmargin=1.6em}
\usepackage{tikz}
\usetikzlibrary{arrows.meta,positioning,fit,calc,backgrounds,shapes.geometric}
\usepackage{algorithm}
\usepackage[noend]{algpseudocode}
\usepackage{titlesec}
\titleformat{\section}{\large\sffamily\bfseries}{\color{accent}\thesection}{0.8em}{}
\titleformat{\subsection}{\normalsize\sffamily\bfseries}{\color{accent}\thesubsection}{0.7em}{}
\titleformat{\paragraph}[runin]{\normalsize\sffamily\bfseries}{}{}{}[.]
\titlespacing*{\section}{0pt}{14pt plus 3pt}{6pt}
\titlespacing*{\subsection}{0pt}{10pt plus 2pt}{4pt}
\titlespacing*{\paragraph}{0pt}{6pt plus 1pt}{0.6em}
\usepackage{fancyhdr}
\pagestyle{fancy}\fancyhf{}
\renewcommand{\headrulewidth}{0pt}
\fancyhead[L]{\footnotesize\sffamily\color{ink2}RamFlow}
\fancyhead[R]{\footnotesize\sffamily\color{ink2}\shorttitle}
\fancyfoot[C]{\footnotesize\sffamily\color{ink2}\thepage}
\fancypagestyle{plain}{\fancyhf{}\fancyfoot[C]{\footnotesize\sffamily\color{ink2}\thepage}}
\usepackage[hidelinks,unicode]{hyperref}
\hypersetup{colorlinks=true,linkcolor=accent,citecolor=accent,urlcolor=accent}
\usepackage{xurl}
\usepackage[numbers,sort&compress]{natbib}
\setlength{\bibsep}{2pt plus 1pt}
\renewcommand{\bibfont}{\small}
\tikzset{prose/.style={}}
\ifko\tikzset{prose/.style={execute at begin picture={\XeTeXlinebreaklocale "ko"\XeTeXlinebreakskip=0pt plus 0.6pt\relax}}}\fi
\newcommand{\KoEn}[2]{\ifko #1\else #2\fi}
\newcommand{\code}[1]{\texttt{#1}}
\newcommand{\pct}{\%}
% Generated by paper/analysis/analyze.py -- do not edit by hand.
\newcommand{\numCSOff}{88.15}
\newcommand{\numCSOn}{48.14}
\newcommand{\numCSRed}{45.4}
\newcommand{\numCSBlkMed}{41.0}
\newcommand{\numCSBlkMin}{39.2}
\newcommand{\numCSBlkMax}{57.3}
\newcommand{\numCSGmRed}{46.5}
\newcommand{\numCSGmLo}{34.9}
\newcommand{\numCSGmHi}{55.9}
\newcommand{\numCSOffFifty}{76.55}
\newcommand{\numCSOnFifty}{38.73}
\newcommand{\numCOOff}{88.56}
\newcommand{\numCOOn}{48.95}
\newcommand{\numCORed}{44.7}
\newcommand{\numCOBlkMed}{39.6}
\newcommand{\numCOBlkMin}{37.2}
\newcommand{\numCOBlkMax}{53.2}
\newcommand{\numCOGmRed}{43.3}
\newcommand{\numCOGmLo}{33.6}
\newcommand{\numCOGmHi}{51.5}
\newcommand{\numCOOffFifty}{76.44}
\newcommand{\numCOOnFifty}{39.35}
\newcommand{\numVCOff}{99.36}
\newcommand{\numVCOn}{50.52}
\newcommand{\numVCRed}{49.2}
\newcommand{\numVCBlkMed}{45.4}
\newcommand{\numVCBlkMin}{43.7}
\newcommand{\numVCBlkMax}{58.1}
\newcommand{\numVCGmRed}{50.5}
\newcommand{\numVCGmLo}{40.4}
\newcommand{\numVCGmHi}{58.9}
\newcommand{\numVCOffFifty}{65.14}
\newcommand{\numVCOnFifty}{38.80}
\newcommand{\numVSOff}{62.46}
\newcommand{\numVSOn}{31.60}
\newcommand{\numVSRed}{49.4}
\newcommand{\numVSBlkMed}{35.8}
\newcommand{\numVSBlkMin}{26.3}
\newcommand{\numVSBlkMax}{63.7}
\newcommand{\numVSGmRed}{43.4}
\newcommand{\numVSGmLo}{19.1}
\newcommand{\numVSGmHi}{60.4}
\newcommand{\numVSOffFifty}{34.97}
\newcommand{\numVSOnFifty}{23.50}
\newcommand{\numBgCOff}{267,405}
\newcommand{\numBgCOn}{107,056}
\newcommand{\numBgCRed}{60.0}
\newcommand{\numBgVOff}{220,064}
\newcommand{\numBgVOn}{87,371}
\newcommand{\numBgVRed}{60.3}
\newcommand{\numEfSOff}{72.70}
\newcommand{\numEfSOn}{44.40}
\newcommand{\numEfOOff}{83.00}
\newcommand{\numEfOOn}{46.90}
\newcommand{\numSessions}{10}
\newcommand{\numSessionsOk}{10}
\newcommand{\numPooledP}{0.0020}
\newcommand{\numTasks}{1,200}
\newcommand{\numConditions}{20}
\newcommand{\numMxOff}{232.85}
\newcommand{\numMxOn}{131.32}
\newcommand{\numMxRed}{43.6}
\newcommand{\numMxBlkMin}{39.7}
\newcommand{\numMxBlkMax}{43.8}
\newcommand{\numMxGmLo}{40.3}
\newcommand{\numMxGmHi}{44.6}
\newcommand{\numMxOffFifty}{223.60}
\newcommand{\numMxOnFifty}{129.01}
\newcommand{\numMxFiftyRed}{42.3}
\newcommand{\numMxBgRed}{59.6}
\newcommand{\numMxAvLo}{856.8}
\newcommand{\numMxAvHi}{1144.0}
\newcommand{\numCpOff}{230.78}
\newcommand{\numCpOn}{219.76}
\newcommand{\numCpRed}{4.8}
\newcommand{\numCpBlkMin}{0.9}
\newcommand{\numCpBlkMax}{6.0}
\newcommand{\numCpGmLo}{1.8}
\newcommand{\numCpGmHi}{6.6}
\newcommand{\numCpOffFifty}{227.62}
\newcommand{\numCpOnFifty}{128.99}
\newcommand{\numCpFiftyRed}{43.3}
\newcommand{\numCpBgRed}{38.4}
\newcommand{\numCpAvLo}{14425.9}
\newcommand{\numCpAvHi}{15157.5}
\newcommand{\numScAmb}{1,811.26}
\newcommand{\numScAbs}{1,896.05}
\newcommand{\numScDry}{1,853.10}
\newcommand{\numScLive}{1,877.65}
\newcommand{\numScDiff}{−18.40}
\newcommand{\numScDiffPct}{−0.97}
\newcommand{\numCiLevel}{93.8}
\newcommand{\numSignP}{0.0312}
\newcommand{\numWilP}{0.0625}

% A rounded "finding" box used for key takeaways
\newcommand{\finding}[1]{%
  \par\smallskip\noindent
  \begin{tikzpicture}[prose]
    \node[fill=panel,rounded corners=3pt,inner sep=7pt,text width=\dimexpr\linewidth-14pt\relax,align=left]{\small #1};
    \fill[accent,rounded corners=1pt] ($(current bounding box.north west)+(0,-1pt)$) rectangle ($(current bounding box.south west)+(2.2pt,1pt)$);
  \end{tikzpicture}\par\smallskip}
